\documentclass[10pt,a4paper]{amsart}
\usepackage[margin=19mm]{geometry}
\usepackage{amsmath,amssymb,mathtools}
\usepackage[scr=rsfs,cal=euler]{mathalfa}
\usepackage{xfrac}
\usepackage[unicode,hidelinks,bookmarks=false]{hyperref}
\newcommand{\ie}{i.e.\ }
\newcommand{\cf}{cf.\ }
\newcommand{\resp}{respectively\ }
\newcommand{\Subsection}{\subsection}
\providecommand{\nobreakdash}{\nobreak\hbox{-}}
\setlength{\emergencystretch}{2em}
\multlinegap0pt
\usepackage{amssymb}
\usepackage{dsfont}
\usepackage{mathtools}
\usepackage[normalem]{ulem}
\usepackage{algorithm}\usepackage{algpseudocode}
\newtheorem{proposition}{Proposition}
\title{\LARGE \bf Linear Systems as Representations of Time Groups}
\author{Subhrajit Sinha}
\date{Source: arXiv:2604.08940v1 (CC BY 4.0)}
\begin{document}
\maketitle

\noindent\emph{Source and licence.} \LARGE \bf Linear Systems as Representations of Time Groups. Version arXiv:2604.08940v1 (CC BY 4.0), distributed by arXiv under the Creative Commons Attribution 4.0 International licence, http://creativecommons.org/licenses/by/4.0/, as recorded in the arXiv licence metadata for this exact version and shown on the abstract page https://arxiv.org/abs/2604.08940v1. Full licence notice: \texttt{licenses/arxiv-2604.08940-CC-BY-4.0.txt}.\par\smallskip

\noindent\textit{Editorial scope.} Counted unit: the article's proposition prop\_group\_algebra\_module (source0.tex lines 244--252). The article's complete original proof of it is delivered with it (source0.tex lines 254--318, one \texttt{proof} environment of 65 lines). No mathematics is written, repaired or re-derived: the statement and the proof are the article's own lines, in the article's own notation, and no retained line is substituted or re-flowed.  No further source-local context is needed: every cross-reference of every retained line resolves inside the two delivered spans.  Nothing was omitted from any retained span, so the package carries no omission record.  No retained line contains a citation, so the wrapper carries no bibliography and no bibliography entry.  No retained line contains a figure environment, an image-inclusion command or drawing code, so no image is delivered and the package declares no asset dependency.  Editorial formatting only, in addition: an AMS \texttt{amsart} preamble that restates the article's own declarations and macros used by the retained lines, namely the environments \texttt{proposition} on separate counters, together with the standard packages those lines need.  The retained environments print in this document as Proposition 1 in order of appearance, whereas the article prints the same items with the numbering of its own full document.  The complete native source of the article as distributed in the arXiv e-print for this exact version is bundled unchanged as \texttt{sources/198180/source0.tex}.
\begin{proposition}\label{prop_group_algebra_module}
Let $G$ be a group, let $k$ be a field, and let $V$ be a vector space over $k$.
Then a linear representation
\begin{eqnarray}
\rho:G \to GL(V)
\end{eqnarray}
is equivalent to giving $V$ the structure of a left $k[G]$-module, where
$k[G]$ denotes the group algebra of $G$ over $k$.
\end{proposition}

\begin{proof}
Suppose first that
\begin{eqnarray}
\rho:G \to GL(V)
\end{eqnarray}
is a linear representation of $G$ on $V$. Recall that the group algebra
$k[G]$ consists of all finite formal sums
\begin{eqnarray}
\sum_{g \in G} a_g g, \qquad a_g \in k,
\end{eqnarray}
with multiplication induced by the group multiplication in $G$ and extended
linearly.

Define an action of $k[G]$ on $V$ by
\begin{eqnarray}
\left(\sum_{g \in G} a_g g\right)\cdot v
=
\sum_{g \in G} a_g \rho(g)v,
\qquad v \in V.
\end{eqnarray}
This is well-defined and linear in both the coefficients and the vector.
Moreover, if
\begin{eqnarray}
\alpha = \sum_{g \in G} a_g g, \qquad
\beta = \sum_{h \in G} b_h h,
\end{eqnarray}
then for any $v\in V$,
\begin{eqnarray*}
(\alpha\beta)\cdot v
&=&
\left(\sum_{g,h \in G} a_g b_h (gh)\right)\cdot v \\
&=&
\sum_{g,h \in G} a_g b_h \rho(gh)v \\
&=&
\sum_{g,h \in G} a_g b_h \rho(g)\rho(h)v \\
&=&
\alpha \cdot (\beta \cdot v).
\end{eqnarray*}
Also, if $e$ is the identity element of $G$, then
\begin{eqnarray*}
e \cdot v = \rho(e)v = Iv = v.
\end{eqnarray*}
Hence $V$ is a left $k[G]$-module.

Conversely, suppose $V$ is a left $k[G]$-module. For each $g \in G$, define
\begin{eqnarray}
\rho(g):V \to V, \qquad \rho(g)v = g \cdot v.
\end{eqnarray}
Each $\rho(g)$ is linear, and for $g,h \in G$ and $v \in V$,
\begin{eqnarray*}
\rho(gh)v = (gh)\cdot v = g \cdot (h \cdot v) = \rho(g)\rho(h)v.
\end{eqnarray*}
Thus,
\begin{eqnarray*}
\rho(gh)=\rho(g)\rho(h), \qquad \rho(e)=I,
\end{eqnarray*}
and each $\rho(g)$ is invertible with inverse $\rho(g^{-1})$. Therefore,
\begin{eqnarray}
\rho:G \to GL(V)
\end{eqnarray}
is a linear representation of $G$ on $V$.

Thus, a linear representation of $G$ on $V$ is equivalent to a left
$k[G]$-module structure on $V$.
\end{proof}

\end{document}
